\section{Introduction}

The preceding hard-window theorem transports an already defined coefficient
covariance to the physical interval with a relative Gram error.  The next
question is therefore not another unsigned energy estimate: it is whether the
literal Möbius--log cluster coefficient can supply a signed principal term in
the two polarized coefficient lanes.

The answer depends on where the multiplier is placed.  If the same scalar
$C_h$ multiplies both lanes in one orthogonal denominator block, conjugate
linearity forces the product $\overline{C_h}C_h=|C_h|^2$.  The aggregate outer
sign or phase disappears exactly.  This does not make the arithmetic tail
unsigned: the Möbius signs inside the sum defining $C_h$ can still alter its
magnitude.  It only rules out a second, external sign mechanism.

We prove three statements.  First, common blockwise phases preserve the exact
coefficient covariance and both norms.  Second, for arbitrary linear
embeddings into a common ambient Hilbert space, the complete dependence on real
sign flips is a finite Walsh polynomial indexed by the edges of the block
graph.  Third, when the direct-sum coefficients enter the hard-window synthesis
of TPC-243, all residual dependence on a common sign pattern is bounded by two
copies of the Gram error.

The contribution is deliberately typed as structural.  Current sources retain
one literal $C_h$ in each packet and canonical primitive frequency coordinates,
but they do not yet identify the literal V59 source sequences with two such
coefficient lanes in one common synthesis map.  No arithmetic cancellation,
coefficient norm bound, or twin-prime conclusion is claimed.
